\subsection{Resumen}
\label{subsec:resumen}

% ####################################################################################################################
% ####################################################################################################################

\textbf{Práctica 1 -- Determinación de la densidad:} Se determinó la densidad de un cilindro macizo mediante tres instrumentos, obteniéndose:

\begin{itemize}
    \item $\rho_{\text{calibre}} = 8.34 \pm 0,065 \, \frac{g}{cm^3}$
    \item $\rho_{\text{regla}} = 8,36 \pm 0,75 \, \frac{g}{cm^3}$
    \item $\rho_{\text{probeta}} = 8,31 \pm 0,91 \, \frac{g}{cm^3}$
\end{itemize}

Los resultados son compatibles entre sí y consistentes con la densidad de la constante de las pesas de latón $\rho = 8,4 \, \frac{g}{cm^3}$.

% ####################################################################################################################
% ####################################################################################################################

\textbf{Práctica 2 -- Péndulo simple:} Se midió el período $T$ para cuatro longitudes $L$ y cuatro amplitudes, verificando la independencia de $T$ respecto a la amplitud en ángulos pequeños.
Además, mediante una regresión lineal de $T^2$ en función de $L$, se determinó el valor de la aceleración de la gravedad:

\begin{equation*}
    g = 9{,}579 \pm 0.255 \frac{m}{s^2}
\end{equation*}
